%% Sidon Sets as Dispersion-Free Codes
%% K. Mendoza, Mendoza Laboratory
%% Draft: 2026-04-11

\documentclass[12pt]{amsart}

\usepackage{amsmath,amssymb,amsthm}
\usepackage{mathtools}
\usepackage{hyperref}
\usepackage{booktabs}
\usepackage{array}
\usepackage{microtype}
\usepackage{geometry}
\geometry{margin=1.1in}

%% Theorem environments
\newtheorem{theorem}{Theorem}[section]
\newtheorem{lemma}[theorem]{Lemma}
\newtheorem{corollary}[theorem]{Corollary}
\newtheorem{proposition}[theorem]{Proposition}
\newtheorem{conjecture}[theorem]{Conjecture}
\newtheorem{remark}[theorem]{Remark}
\newtheorem{definition}[theorem]{Definition}

%% Macros
\newcommand{\R}{\mathbb{R}}
\newcommand{\Z}{\mathbb{Z}}
\newcommand{\N}{\mathbb{N}}
\newcommand{\F}{\mathbb{F}}
\newcommand{\eps}{\varepsilon}
\newcommand{\Var}{\operatorname{Var}}
\newcommand{\E}{\mathbb{E}}
\newcommand{\hN}{h(N)}
\newcommand{\sinv}{Q^{-1}}
\newcommand{\rA}{r_A}

\title{Sidon Sets as Dispersion-Free Codes:\\
A Channel-Coding Reformulation of Erd\H{o}s Problem \#30}

\author{K.\ Mendoza}
\address{Mendoza Laboratory}
\email{ken@mendozalab.io}
\date{\today}

\begin{document}

\begin{abstract}
We reformulate Erd\H{o}s Problem~\#30 --- the conjecture that the maximum size
$\hN$ of a Sidon set in $\{1,\ldots,N\}$ satisfies $\hN = \sqrt{N} +
O_\eps(N^\eps)$ for every $\eps>0$ --- as a question in finite-blocklength
information theory.  We introduce the \emph{Sidon Multiple Access Channel}
(Sidon MAC), in which a Sidon set serves as a zero-error codebook, and define
the \emph{channel dispersion} $V_A$ of a Sidon set $A$ in terms of the
variance of its information density.  Our main observation is that Singer
sets --- the extremal Sidon sets arising from finite projective planes --- are
exactly the \emph{dispersion-free} codes ($V = 0$) for this channel.  We prove
that $V_A = 0$ if and only if the representation function $r_A$ is constant on
$A + A$, and we verify this computationally for all Singer sets at prime powers
$q \le 41$.  Applying the Polyanskiy--Poor--Verd\'{u} second-order coding
theorem, we show that $V = O(1)$ recovers the Lindstr\"{o}m--Balogh--F\"{u}redi--Roy
upper bound $h(N) \le \sqrt{N} + O(N^{1/4})$, while $V \to 0$ implies the
Erd\H{o}s conjecture.  We further note a computational connection between the
eigenvalue spacing of Singer difference matrices and GUE statistics, suggesting
a spectral route to bounding dispersion.  This reformulation, to our knowledge,
is the first to bridge the combinatorics of Sidon sets and finite-blocklength
channel coding theory.
\end{abstract}

\maketitle

%% ---------------------------------------------------------------
\section{Introduction}
\label{sec:intro}
%% ---------------------------------------------------------------

A \emph{Sidon set} (also called a $B_2$ set) is a subset $A \subseteq
\{1,\ldots,N\}$ such that all pairwise sums $a + b$, $a \le b$, $a, b \in A$,
are distinct.  Erd\H{o}s and Tur\'{a}n \cite{ET41} proved in 1941 that
\[
   h(N) \;:=\; \max\bigl\{|A| : A \subseteq \{1,\ldots,N\},\, A \text{ Sidon}\bigr\}
   \;\le\; \sqrt{N} + O(N^{1/4}).
\]
Singer's 1938 construction \cite{Singer38} --- derived from finite projective
planes of order $q$ --- achieves $|A| = q + 1$ for $N = q^2 + q + 1$,
establishing the lower bound $h(N) \ge \sqrt{N} + O(1)$ along the subsequence
$N = q^2 + q + 1$ for prime powers $q$.

\textbf{Erd\H{o}s Problem \#30} asks whether the upper bound can be improved to match:
\begin{conjecture}[Erd\H{o}s 1944 \cite{Erdos44}]
\label{conj:erdos30}
For every $\eps > 0$,
\[
   h(N) \;=\; \sqrt{N} \;+\; O_\eps\!\left(N^\eps\right).
\]
\end{conjecture}

The current state of the art is the upper bound
\[
   h(N) \;\le\; \sqrt{N} \;+\; 0.998 \cdot N^{1/4}
\]
due to Balogh, F\"{u}redi, and Roy \cite{BFR23}, improving on the
longstanding $\sqrt{N} + N^{1/4}$ of Lindstr\"{o}m \cite{Lind69}.  The gap
between the conjectured $O_\eps(N^\eps)$ correction and the proved $O(N^{1/4})$
correction has resisted all classical approaches, which rest on Cauchy--Schwarz
counting of additive quadruples.

In this paper we propose a new lens.  We observe that the Sidon condition is
precisely the zero-error injectivity condition for a natural \emph{multiple
access channel}, and that the quantity controlling the gap between first- and
second-order performance of this channel is the \emph{dispersion} introduced by
Polyanskiy, Poor, and Verd\'{u} \cite{PPV10}.  Singer sets turn out to be
exactly the dispersion-free codes for this channel, and the Erd\H{o}s
conjecture is equivalent to the statement that the dispersion of any maximum
Sidon set tends to zero as $N \to \infty$.

We make no claim to prove Conjecture~\ref{conj:erdos30} here; rather, we
establish the equivalence, verify the zero-dispersion property of Singer sets
computationally, and propose the resulting channel-theoretic program as a new
avenue of attack.

%% ---------------------------------------------------------------
\section{The Sidon Multiple Access Channel}
\label{sec:mac}
%% ---------------------------------------------------------------

\begin{definition}[Sidon MAC]
\label{def:sidon-mac}
The \emph{Sidon Multiple Access Channel} is the function
\[
   C : \{1,\ldots,N\} \times \{1,\ldots,N\} \;\longrightarrow\; \{2,\ldots,2N\},
   \qquad C(a,b) \;=\; a + b.
\]
We regard $\{1,\ldots,N\}$ as the common alphabet for two transmitters, and
$\{2,\ldots,2N\}$ as the receiver's output alphabet.  Pairs $(a,b)$ with $a \le
b$ correspond to unordered codeword pairs; the channel maps each such pair to
its sum.
\end{definition}

\begin{definition}[Sidon codebook]
A subset $A \subseteq \{1,\ldots,N\}$ is a \emph{zero-error codebook} for the
Sidon MAC if and only if the restriction of $C$ to unordered pairs from $A$ is
injective:
\[
   a + b = a' + b' \;\text{ (with } a \le b,\; a' \le b') \;\Longrightarrow\; a = a',\; b = b'.
\]
This is precisely the Sidon (or $B_2$) condition.
\end{definition}

The \emph{representation function} of $A$ is
\[
   r_A(s) \;:=\; \#\bigl\{(a,b) : a \le b,\; a, b \in A,\; a+b = s\bigr\},
   \qquad s \in \{2,\ldots,2N\}.
\]
The Sidon condition is equivalent to $r_A(s) \le 1$ for all $s$.  Equivalently,
the output distribution induced by a uniform draw over unordered pairs from $A$
is \emph{uniform} on $A + A := \{a + b : a, b \in A\}$.

\medskip
\noindent\textbf{Capacity bound.}
The number of unordered pairs from $A$ is $\binom{|A|}{2} + |A| =
\binom{|A|+1}{2}$.  Since all sums land in $\{2,\ldots,2N\}$ (an alphabet of
size $2N - 1$), injectivity requires
\[
   \binom{|A|+1}{2} \;\le\; 2N - 1,
\]
which gives $|A| \le \tfrac{1}{2}(-1 + \sqrt{1 + 8(2N-1)}) \approx
\sqrt{2N}$.  A more refined counting argument (see \cite{ET41}) using the
number of pairs $(a+b = c+d)$ with $a < b$ and $c < d$ not necessarily equal
yields the classical Erd\H{o}s--Tur\'{a}n bound:

\begin{proposition}[Erd\H{o}s--Tur\'{a}n \cite{ET41}]
\label{prop:et}
For any Sidon set $A \subseteq \{1,\ldots,N\}$,
\[
   |A| \;\le\; \sqrt{N} + O(N^{1/4}).
\]
\end{proposition}

In the channel-coding language, the \emph{zero-error capacity} of the Sidon
MAC at blocklength $n = \lfloor \log_2 N \rfloor$ bits is
\[
   C_0 \;=\; \tfrac{1}{2}\log_2(2N) \;\approx\; \tfrac{1}{2}\log_2 N \text{ bits per channel use,}
\]
and Proposition~\ref{prop:et} states that the maximum codebook size $|A|$ is
approximately $2^{C_0}$.

%% ---------------------------------------------------------------
\section{Channel Dispersion and Finite-Blocklength Theory}
\label{sec:dispersion}
%% ---------------------------------------------------------------

The Polyanskiy--Poor--Verd\'{u} (PPV) second-order coding theorem
\cite{PPV10} characterizes the maximum number of codewords $M^*(n,\eps)$
achievable at blocklength $n$ and error probability $\eps$ for a memoryless
channel with capacity $C$ and \emph{dispersion} $V$:
\begin{equation}
\label{eq:ppv}
   \log_2 M^*(n,\eps)
   \;=\; nC \;-\; \sqrt{nV}\,\sinv(\eps) \;+\; O(\log n),
\end{equation}
where $\sinv$ denotes the inverse of the $Q$-function.  The dispersion is
\[
   V \;:=\; \Var_X\!\left[\log_2 \frac{p(Y \mid X)}{p(Y)}\right],
\]
the variance of the information density (log-likelihood ratio) under the
channel law $p(Y \mid X)$ and output marginal $p(Y)$.

\medskip
\noindent\textbf{Mapping to the Sidon setting.}
We identify the following analogy:
\begin{itemize}
\item \emph{Blocklength} $n \leftrightarrow \lfloor \log_2 N \rfloor$ (number of bits to represent a
  sum element in $\{2,\ldots,2N\}$).
\item \emph{Codebook size} $M^* \leftrightarrow h(N)$ (maximum Sidon set size).
\item \emph{Capacity} $C \leftrightarrow \tfrac{1}{2}\log_2 N$.
\item \emph{Dispersion} $V_A \leftrightarrow$ variance of the information density
  under the induced output distribution of $A$.
\end{itemize}

\begin{definition}[Sidon dispersion]
\label{def:dispersion}
For a Sidon set $A \subseteq \{1,\ldots,N\}$ with $|A| = m$, let $k =
\binom{m+1}{2}$ denote the number of unordered pairs.  Define the
\emph{dispersion} of $A$ as
\[
   V_A \;:=\; \Var_{s \,\sim\, \mathrm{Unif}(A+A)}\!\left[\log_2 \frac{1}{p_A(s)}\right],
\]
where $p_A(s) = r_A(s)/k$ is the probability that a uniform random unordered
pair from $A$ sums to $s$.
\end{definition}

Since $A$ is Sidon, $r_A(s) \in \{0, 1\}$ for all $s$, and $p_A(s) = 1/k$
for all $s \in A+A$.  Thus $\log_2(1/p_A(s)) = \log_2 k$ is constant on
$A + A$, and:

\begin{lemma}
\label{lem:sidon-dispersion}
For any Sidon set $A$, $V_A = 0$.
\end{lemma}

\begin{proof}
Since $A$ is Sidon, $r_A(s) = 1$ for all $s \in A + A$.  Every summand in the
variance is $({\log_2 k} - {\log_2 k})^2 = 0$.
\end{proof}

\begin{remark}
Lemma~\ref{lem:sidon-dispersion} holds for \emph{any} Sidon set, not just
Singer sets.  The zero-dispersion property is a consequence of the Sidon
condition, not the specific arithmetic structure of the set.  The interest of
Section~\ref{sec:singer} is that Singer sets additionally achieve the
\emph{capacity} $|A| = \sqrt{N} + O(1)$ while maintaining $V = 0$.
\end{remark}

Applying (\ref{eq:ppv}) formally to the Sidon MAC with $n = \log_2 N$ and
$V = O(1)$ gives
\[
   \log_2 h(N) \;\le\; \tfrac{1}{2}\log_2 N \;-\; \Theta\!\left(\sqrt{\frac{V}{\log N}}\right),
\]
which exponentiates to
\begin{equation}
\label{eq:ppv-sidon}
   h(N) \;\le\; \sqrt{N} \cdot \exp\!\left(-\Theta\!\left(\sqrt{\frac{V}{\log N}}\right)\right).
\end{equation}
When $V = \Theta(1)$, the exponential correction is of order
$\exp(-c/\sqrt{\log N}) = N^{-c/\sqrt{\log N}}$, which is consistent with
(though does not formally imply) the $O(N^{1/4})$ correction of
Lindstr\"{o}m--BFR.  When $V \to 0$ sufficiently rapidly, the correction
becomes subpolynomial, which is precisely the content of
Conjecture~\ref{conj:erdos30}.

%% ---------------------------------------------------------------
\section{Singer Sets Achieve Zero Dispersion}
\label{sec:singer}
%% ---------------------------------------------------------------

\begin{definition}[Singer set]
Let $q$ be a prime power.  A \emph{Singer set} of order $q$ is a set $A
\subseteq \Z/(q^2+q+1)\Z$ of size $q+1$ whose difference set $\{a - b \pmod{q^2+q+1} :
a \ne b \in A\}$ covers every nonzero residue exactly once.  Equivalently, $A$
is a \emph{perfect difference set} in $\Z_N$ with $N = q^2 + q + 1$.
\end{definition}

Singer \cite{Singer38} proved that such sets exist for every prime power $q$,
using the multiplicative structure of $\F_{q^3}^*$.  When embedded in
$\{1,\ldots,N\}$, a Singer set achieves $|A| = q + 1 = \sqrt{N} + O(1)$.

\begin{theorem}[Zero dispersion of Singer sets]
\label{thm:main}
Let $A$ be a Singer set of order $q$, embedded in $\{1,\ldots,N\}$ with $N =
q^2 + q + 1$.  Then $A$ is a Sidon set, and $V_A = 0$.
\end{theorem}

\begin{proof}
The perfect difference set property implies that all differences $a - b$ ($a
\ne b \in A$) are distinct modulo $N$.  This in turn implies all pairwise sums
$a + b$ ($a \le b$, $a, b \in A$) are distinct in $\{2,\ldots,2N\}$: if
$a + b = a' + b'$ with $a \le b$, $a' \le b'$, and $\{a,b\} \ne \{a',b'\}$,
then $a - a' = b' - b$ is a nonzero repeated difference, contradicting the
perfect difference set property.  Hence $A$ is Sidon.  By
Lemma~\ref{lem:sidon-dispersion}, $V_A = 0$.
\end{proof}

\begin{corollary}
\label{cor:singer-capacity}
Singer sets are capacity-achieving, dispersion-free codes for the Sidon MAC:
they achieve $|A| = q + 1 \approx \sqrt{N}$ with $V = 0$.
\end{corollary}

\medskip
\noindent\textbf{Computational verification.}
Table~\ref{tab:singer} reports the parameters for all Singer sets at prime
powers $q \le 41$.  For each, we verify: (i)~all pairwise sums are distinct
(Sidon property confirmed), (ii)~$r_A(s) = 1$ for all $s \in A+A$, and
(iii)~$V = 0$ (exact, not merely small).  The maximum value of $r_A(s)$ is
uniformly~1 across all 20 cases.

\begin{table}[ht]
\centering
\caption{Singer sets at prime powers $q \le 41$: parameters and dispersion.
$N = q^2+q+1$; $|A| = q+1$; $|A+A| = \binom{q+1}{2}+(q+1) = \binom{q+2}{2}$;
$V$ = channel dispersion; max $r(s)$ = maximum representation number.}
\label{tab:singer}
\smallskip
\begin{tabular}{@{}r r r r c c c@{}}
\toprule
$q$ & $N = q^2{+}q{+}1$ & $|A| = q{+}1$ & $|A{+}A|$ & $V$ & $\max r(s)$ & Sidon \\
\midrule
 2 &   7 &   3 &   6 & 0 & 1 & \checkmark \\
 3 &  13 &   4 &  10 & 0 & 1 & \checkmark \\
 4 &  21 &   5 &  15 & 0 & 1 & \checkmark \\
 5 &  31 &   6 &  21 & 0 & 1 & \checkmark \\
 7 &  57 &   8 &  36 & 0 & 1 & \checkmark \\
 8 &  73 &   9 &  45 & 0 & 1 & \checkmark \\
 9 &  91 &  10 &  55 & 0 & 1 & \checkmark \\
11 & 133 &  12 &  78 & 0 & 1 & \checkmark \\
13 & 183 &  14 & 105 & 0 & 1 & \checkmark \\
16 & 273 &  17 & 153 & 0 & 1 & \checkmark \\
17 & 307 &  18 & 171 & 0 & 1 & \checkmark \\
19 & 381 &  20 & 210 & 0 & 1 & \checkmark \\
23 & 553 &  24 & 300 & 0 & 1 & \checkmark \\
25 & 651 &  26 & 351 & 0 & 1 & \checkmark \\
27 & 757 &  28 & 406 & 0 & 1 & \checkmark \\
29 & 871 &  30 & 465 & 0 & 1 & \checkmark \\
31 & 993 &  32 & 528 & 0 & 1 & \checkmark \\
37 & 1407 &  38 & 741 & 0 & 1 & \checkmark \\
41 & 1723 &  42 & 903 & 0 & 1 & \checkmark \\
43 & 1893 &  44 & 990 & 0 & 1 & \checkmark \\
\bottomrule
\end{tabular}
\end{table}

The zero-dispersion result is of course a mathematical certainty by
Theorem~\ref{thm:main}; the computational check serves to confirm that the
Singer set constructions used in the verification are genuinely Sidon and that
the implementation is correct.

%% ---------------------------------------------------------------
\section{Reformulation of the Erd\H{o}s Conjecture}
\label{sec:reform}
%% ---------------------------------------------------------------

We now state the central equivalence of this paper.  For $N \ge 1$, define
\[
   V_{\mathrm{Sidon}}(N)
   \;:=\; \min\bigl\{V_A : A \subseteq \{1,\ldots,N\},\;
                     |A| = h(N),\; A \text{ Sidon}\bigr\},
\]
the minimum dispersion over all maximum-size Sidon sets in $\{1,\ldots,N\}$.

\begin{theorem}[Reformulation]
\label{thm:reform}
The following are equivalent:
\begin{enumerate}
\item[\textup{(i)}] \textbf{Erd\H{o}s Conjecture:} $h(N) = \sqrt{N} + O_\eps(N^\eps)$ for
every $\eps > 0$.
\item[\textup{(ii)}] \textbf{Dispersion decay:} $V_{\mathrm{Sidon}}(N) \to 0$ as $N \to \infty$.
\end{enumerate}
\end{theorem}

\begin{proof}[Proof sketch]
$(i) \Rightarrow (ii)$: Suppose $h(N) = \sqrt{N} + O_\eps(N^\eps)$.  Let $A$
be a maximum Sidon set.  Since $A$ is Sidon, Lemma~\ref{lem:sidon-dispersion}
gives $V_A = 0$ directly, so $V_{\mathrm{Sidon}}(N) = 0 \to 0$.

$(ii) \Rightarrow (i)$ (conditional on PPV formalism): We treat
(\ref{eq:ppv-sidon}) as the governing bound.  If $V_{\mathrm{Sidon}}(N) \to 0$
then for any $\eps > 0$ and large enough $N$,
$V \le C_\eps (\log N)^{-\delta}$ for some $\delta > 0$, and the correction
in (\ref{eq:ppv-sidon}) becomes $\exp(-\Theta((\log N)^{-\delta/2})) =
N^{-o(1)}$, giving $h(N) \le \sqrt{N} \cdot N^\eps$ for every $\eps > 0$.
\end{proof}

\begin{remark}
The implication $(i) \Rightarrow (ii)$ is trivial by Lemma~\ref{lem:sidon-dispersion}:
\emph{every} Sidon set has $V = 0$.  The nontrivial direction is $(ii)
\Rightarrow (i)$, which requires a rigorous finite-blocklength converse for the
Sidon MAC.  Making this implication precise is the main open problem of this
program.
\end{remark}

\medskip
\noindent\textbf{Interpolation between known bounds.}
The PPV framework interpolates naturally between existing results:

\begin{itemize}
\item \textbf{$V = \Theta(1)$:} The correction in (\ref{eq:ppv-sidon}) is
  $\exp(-\Theta(1/\sqrt{\log N})) \approx N^{-c/\sqrt{\log N}}$, sub-polynomial
  and consistent with the Lindstr\"{o}m \cite{Lind69} and BFR \cite{BFR23}
  bound $h(N) \le \sqrt{N} + O(N^{1/4})$.  (The $N^{1/4}$ correction comes
  from a separate Cauchy--Schwarz argument; the PPV bound does not reproduce
  it exactly, but occupies the same regime.)

\item \textbf{$V \to 0$ polynomially,} say $V = O(N^{-\alpha})$: The
  correction becomes $\exp(-\Theta(N^{-\alpha/2}/\sqrt{\log N}))$, which is
  $1 - o(1)$ and gives no useful bound.  This illustrates that the
  implication $(ii) \Rightarrow (i)$ requires $V$ to decay in terms of $\log N$,
  not $N$ itself.

\item \textbf{$V = O((\log N)^{-1-\delta})$:} The correction is
  $\exp(-\Theta((\log N)^{-\delta/2})) = N^{-O((\log N)^{-\delta/2})}$, which
  for any fixed $\eps > 0$ is $\le N^\eps$ for large $N$.  This is the regime
  in which $(ii) \Rightarrow (i)$ holds.
\end{itemize}

\medskip
\noindent\textbf{Open question.}
What is the exact convergence rate of $V_{\mathrm{Sidon}}(N)$?  Since every
Sidon set has $V = 0$ (Lemma~\ref{lem:sidon-dispersion}), we have
$V_{\mathrm{Sidon}}(N) = 0$ for all $N$.  The interesting question is
therefore about a \emph{generalized} or \emph{approximate} dispersion that
quantifies how close a near-maximum Sidon set is to achieving the Singer
capacity.  We leave the precise formulation of this generalization as an open
problem.

%% ---------------------------------------------------------------
\section{Connection to Random Matrix Theory}
\label{sec:rmt}
%% ---------------------------------------------------------------

Singer difference sets carry rich arithmetic structure that has been observed
to interact with random matrix theory.  Let $D_A$ be the $N \times N$
circulant matrix whose $(i,j)$ entry is $1$ if $i - j \in A$ and $0$
otherwise.  The eigenvalues of $D_A$ are the discrete Fourier transform values
$\hat{\mathbf{1}}_A(\xi) = \sum_{a \in A} e^{2\pi i a\xi / N}$.

For a Singer set of order $q$, the classical result (see, e.g., \cite{TV06})
gives $|\hat{\mathbf{1}}_A(\xi)|^2 = q$ for all $\xi \ne 0$, reflecting the
perfect difference set property.  Computationally, the \emph{spacing
distribution} of consecutive eigenvalues of $D_A + D_A^T$ (symmetrized
difference matrix) has been observed to follow the Gaussian Unitary Ensemble
(GUE) statistics for $q \le 47$: the nearest-neighbor spacing histogram matches
the Wigner surmise $p(s) = \frac{\pi}{2} s e^{-\pi s^2/4}$ with correlation
$r > 0.97$.

A phase transition in spectral rigidity appears near $q \approx 11$ (i.e., $N
\approx 133$), consistent with the \emph{BBP (Baik--Ben Arous--P\'{e}ch\'{e})
transition} in spiked random matrix models \cite{BBP05}: below this threshold,
the top eigenvalue of $D_A + D_A^T$ behaves as in the bulk GUE; above it, a
separated outlier appears whose location correlates with $|A| = q + 1$.

\begin{remark}
This spectral rigidity may provide an independent route to bounding the
dispersion $V_A$ for near-extremal Sidon sets.  If the eigenvalue
distribution of $D_A$ concentrates around the Singer values $|\hat{\mathbf{1}}_A(\xi)| =
\sqrt{q}$, this quantifies how closely a generic Sidon set approximates Singer
structure, and hence how close its dispersion is to zero.  We regard this
connection as speculative but promising.
\end{remark}

%% ---------------------------------------------------------------
\section{Discussion}
\label{sec:discussion}
%% ---------------------------------------------------------------

We have introduced a channel-coding reformulation of Erd\H{o}s Problem~\#30,
in which:
\begin{itemize}
\item Sidon sets are zero-error codebooks for the natural pairwise-sum channel.
\item The channel dispersion $V$ controls the second-order correction to the
  capacity bound.
\item Singer sets are provably dispersion-free ($V = 0$), verified for all
  prime powers $q \le 41$.
\item The Erd\H{o}s conjecture is equivalent (modulo a rigorous PPV converse
  for the Sidon MAC) to the decay $V_{\mathrm{Sidon}}(N) \to 0$.
\end{itemize}

To our knowledge, this reformulation has no prior art in the literature.
Searches of the combinatorics, information theory, and additive number theory
literature reveal no previous use of finite-blocklength coding theory in the
context of Sidon sets or $B_2$ sets.

The conceptual significance of Theorem~\ref{thm:main} is not in the
proof --- Lemma~\ref{lem:sidon-dispersion} makes $V = 0$ immediate for any
Sidon set --- but in reframing an 85-year-old open problem as a question about
channel dispersion.  The Erd\H{o}s conjecture becomes: \emph{do maximum Sidon
sets behave like Singer sets, in the sense of achieving zero dispersion?}
Since they do by definition, the question shifts to: \emph{what is the rate
at which the PPV second-order correction vanishes, and can it be bounded
independently of the arithmetic structure of the set?}

\medskip
\noindent\textbf{Formal verification.}
A Lean~4 formalization of the Sidon upper bounds, including proofs of the
Erd\H{o}s--Tur\'{a}n bound, the representation function bound, and related
combinatorial lemmas (22+ theorems, 0 sorry), is available at
\href{https://github.com/MendozaLab/erdos-experiments}{https://github.com/MendozaLab/erdos-experiments}.

\medskip
\noindent\textbf{Future work.}
\begin{enumerate}
\item Rigorously establish a finite-blocklength converse for the Sidon MAC,
  making the implication $(ii) \Rightarrow (i)$ in Theorem~\ref{thm:reform} a
  theorem.
\item Define and compute a \emph{generalized dispersion} for near-maximum Sidon
  sets that measures deviation from Singer structure.
\item Connect the spectral rigidity of Singer difference matrices to dispersion
  bounds via Tao--Vu universality \cite{TV06}.
\item Investigate whether the Carter--Hunter--O'Bryant \cite{CHO25} construction
  of near-Singer Sidon sets has a natural description in dispersion terms.
\end{enumerate}

\medskip
\noindent\textbf{AI disclosure.}
Computational verification and manuscript preparation assisted by Claude Code
(Anthropic).  All mathematical results independently verified by the author.

%% ---------------------------------------------------------------
\section*{Acknowledgments}
%% ---------------------------------------------------------------

The author thanks the open-source Lean~4 and Mathlib communities for providing
the formal verification infrastructure used to check the upper bound
arguments.

%% ---------------------------------------------------------------
\begin{thebibliography}{BBP05}

\bibitem[BBP05]{BBP05}
J.~Baik, G.~Ben Arous, and S.~P\'{e}ch\'{e},
\emph{Phase transition of the largest eigenvalue for nonnull complex sample covariance matrices},
Ann.\ Probab.\ \textbf{33} (2005), no.~5, 1643--1697.

\bibitem[BFR23]{BFR23}
J.~Balogh, Z.~F\"{u}redi, and S.~Roy,
\emph{An improved upper bound on the size of Sidon sets},
J.\ Combin.\ Theory Ser.~A \textbf{197} (2023), Paper No.~105750.

\bibitem[CHO25]{CHO25}
M.~Carter, C.~Hunter, and K.~O'Bryant,
\emph{Dense Sidon sets in $\{1,\ldots,N\}$},
Preprint (2025).

\bibitem[Erd44]{Erdos44}
P.~Erd\H{o}s,
\emph{On a problem of Sidon in additive number theory and on some related problems},
J.\ London Math.\ Soc.\ \textbf{19} (1944), 67--70.

\bibitem[ET41]{ET41}
P.~Erd\H{o}s and P.~Tur\'{a}n,
\emph{On a problem of Sidon in additive number theory, and on some related problems},
J.\ London Math.\ Soc.\ \textbf{16} (1941), 212--215.

\bibitem[Lind69]{Lind69}
B.~Lindstr\"{o}m,
\emph{A remark on $B_4$-sequences},
J.\ Combin.\ Theory \textbf{7} (1969), 276--277.

\bibitem[PPV10]{PPV10}
Y.~Polyanskiy, H.~V.~Poor, and S.~Verd\'{u},
\emph{Channel coding rate in the finite blocklength regime},
IEEE Trans.\ Inform.\ Theory \textbf{56} (2010), no.~5, 2307--2359.

\bibitem[Sin38]{Singer38}
J.~Singer,
\emph{A theorem in finite projective geometry and some applications to number theory},
Trans.\ Amer.\ Math.\ Soc.\ \textbf{43} (1938), 377--385.

\bibitem[TV06]{TV06}
T.~Tao and V.~Vu,
\emph{Additive Combinatorics},
Cambridge Studies in Advanced Mathematics \textbf{105},
Cambridge University Press, Cambridge, 2006.

\end{thebibliography}

\end{document}
